\section{Stability from a nearly norming head test}\label{stab:section}

A common head test controls each endpoint of a convex combination,
even when the combination has a nonzero tail.  The two relevant errors
are the evaluation deficit and the tail of the combination.  Keeping
them separate also identifies the cost of unequal endpoint weights.

Throughout this section, $Z$ is either the forest coordinate space
$X$ or the infinite-height coordinate predual $B_\infty$.  We use the
notation and estimates of the common atom toolkit: $H$ is a finite
ancestral head, $G_H$ is its set of exterior gates, and $p=P_Hp\in J$
is a head test with $\rho(p)\le1$.  For a represented atom
$A=\sum_jc_jd_{S_j}$, its replacement is
$A^\circ=\sum_js_jd_{S_j}$, where $s_j=p(S_j)$.  The quantities
$L(A)=1-A(p)$, $D(A)=1-\sum_js_j^2$, and
$E(A)=\sum_j(c_j-s_j)^2$ are respectively its evaluation deficit,
squared-test slack, and squared coefficient error.

\subsection{A common test and weighted endpoint laws}

\begin{theorem}\label{stab:weighted}
Let $z_1,\ldots,z_m\in B_Z$ and let
$\lambda_1,\ldots,\lambda_m>0$ with $\sum_i\lambda_i=1$.  Put
\[
 z=\sum_i\lambda_i z_i,\qquad
 \delta_i=1-z_i(p),\qquad
 \delta=1-z(p)=\sum_i\lambda_i\delta_i,
 \qquad \eta=\|Q_Hz\|.
\]
Then, for each $i$,
\begin{equation}\label{stab:weighted-bound}
 \|Q_Hz_i\|
 \le\lambda_i^{-1}
       \sqrt{(\eta+\sqrt{2\delta})^2+4\delta}
       +\sqrt{2\delta_i}.
\end{equation}
All deficits are nonnegative, and $\delta_i\le\delta/\lambda_i$.
For two equally weighted endpoints this gives
\begin{equation}\label{stab:nonzero-midpoint}
 \|Q_Hz_\pm\|
 \le2\sqrt{(\eta+\sqrt{2\delta})^2+4\delta}+2\sqrt\delta.
\end{equation}
\end{theorem}

\begin{proof}
First suppose each endpoint has a finite law of represented atoms.
Write $\mathbb E_i$ for its law and
$\mathbb E=\sum_i\lambda_i\mathbb E_i$ for their mixture.
The pointwise budget of Lemma~\ref{stab:coef:pointwise} gives
\begin{equation}\label{stab:conditional-deficits}
 \mathbb E(D+E)\le2\delta,\qquad
 \mathbb E_i(D+E)\le2\delta_i,
 \qquad
 \|\mathbb E(A-A^\circ)\|\le\sqrt{2\delta}.
\end{equation}
The same argument under the $i$th law bounds its replacement error
by $\sqrt{2\delta_i}$.

Let $s_g(A)$ denote the head sum of the crossing at $g$, with value
zero when no segment crosses there.  A segment wholly outside $H$
has replacement coefficient zero.  Thus $Q_HA^\circ$ consists of
segments starting at the gates, with coefficients $s_g(A)$.
Put $v=\mathbb E s$ and $m=\mathbb E|s|$, coordinatewise.
Gate-coordinate contraction and \eqref{stab:conditional-deficits}
give $\|v\|_2\le\eta+\sqrt{2\delta}$.  Applying the
signed--unsigned identity \eqref{stab:signed-unsigned} to the
mixture yields
\[
 \|m\|_2^2\le\|v\|_2^2+2\mathbb ED
 \le(\eta+\sqrt{2\delta})^2+4\delta.
\]
Equivalently, its positive and negative mean masses $l_g,n_g$
satisfy $\sum_gl_gn_g\le\delta$.

The conditional absolute masses $m_i=\mathbb E_i|s|$ satisfy
$m_i\le m/\lambda_i$ coordinatewise.  Rooted-mass recovery,
Lemma~\ref{stab:coef:tail-assembly}, therefore gives
\[
 \|Q_H\mathbb E_iA^\circ\|
 \le\|m_i\|_2\le\lambda_i^{-1}\|m\|_2.
\]
This comparison concerns nonnegative Euclidean mass vectors; it
does not require lattice monotonicity of the coordinate-space norm.
Adding the conditional replacement error proves
\eqref{stab:weighted-bound}.

For completed endpoints, choose finite coordinate approximants and
divide each by the maximum of one and its norm.  They converge in
norm inside $B_Z$ and have finite atom laws by the finite-hull
theorem for the chosen model.  Evaluation and both coordinate
projections are continuous, so the inequality passes to the limit.
Finally, nonnegativity of the individual deficits gives
$\delta_i\le\delta/\lambda_i$ and hence
\eqref{stab:nonzero-midpoint}.
\end{proof}

\begin{corollary}[Bounded multipliers]\label{stab:bounded-multiplier}
Let $A$ be a finite random represented atom, and let $\nu$ be a real
random variable on the same probability space with $|\nu|\le1$.
Set $x=\mathbb EA$, $u=\mathbb E(\nu A)$,
$\delta=1-x(p)$, and $\eta=\|Q_Hx\|$.  Then
\begin{equation}\label{stab:bounded-multiplier-bound}
 \|Q_Hu\|
 \le\sqrt{(\eta+\sqrt{2\delta})^2+4\delta}+\sqrt{2\delta}
 \le\eta+(2+2\sqrt2)\sqrt\delta.
\end{equation}
Consequently, whenever $x,u\in Z$ satisfy $x\pm u\in B_Z$,
\begin{equation}\label{stab:symmetric-tail}
 \|Q_Hu\|\le\|Q_Hx\|+(2+2\sqrt2)\sqrt{1-x(p)}.
\end{equation}
\end{corollary}

\begin{proof}
Apply \eqref{stab:signed-unsigned} to the head-sum arrays of $A$.
The pointwise replacement budget gives
$\|\mathbb Es\|_2\le\eta+\sqrt{2\delta}$, and hence
\[
 \|\mathbb E|s|\|_2
 \le\sqrt{(\eta+\sqrt{2\delta})^2+4\delta}.
\]
The bounded-multiplier case of
Lemma~\ref{stab:coef:tail-assembly}, applied to the replaced atoms
$A^\circ$, bounds $\|Q_H\mathbb E(\nu A^\circ)\|$ by this
absolute mass.  The remaining term satisfies
\[
 \|Q_H\mathbb E(\nu(A-A^\circ))\|
 \le\mathbb E\sqrt{E(A)}\le\sqrt{2\delta}.
\]
Adding proves the first inequality.  The second follows from
$\sqrt{a^2+b^2}\le a+b$ for $a,b\ge0$.  No independence between
$\nu$ and $A$ was used.

For finite convex representations of $x+u$ and $x-u$, take their
equal mixture and let $\nu$ be its endpoint sign.  Then
$\mathbb EA=x$ and $\mathbb E(\nu A)=u$.  For completed endpoints,
approximate both inside their finite atomic balls and pass to the
limit, as in Theorem~\ref{stab:weighted}.
\end{proof}

\begin{corollary}\label{stab:supported-pair}
Let $z_+,z_-\in B_Z$.  If $z=(z_++z_-)/2$ is supported on $H$
and $\|z\|\ge1-\varepsilon$, where $0<\varepsilon<1$, then
\begin{equation}\label{stab:supported-bound}
 \|Q_Hz_\pm\|\le(2+2\sqrt6)\sqrt\varepsilon.
\end{equation}
More generally, let $R>0$, $d\ge0$, and $h_+,h_-\in Z$ satisfy
$\|h_\pm\|\le R$, with $h_++h_-$ supported on $H$.
If a common head test $p$, $\rho(p)=1$, satisfies
$h_\pm(p)\ge(1-d)R$, then
\begin{equation}\label{stab:common-test-bound}
 \|Q_Hh_\pm\|\le R(2\sqrt{6d}+\sqrt{2d}).
\end{equation}
\end{corollary}

\begin{proof}
For the first assertion choose a norming test in the finite-dimensional
head.  Then $\delta\le\varepsilon$ and $\eta=0$, so
\eqref{stab:nonzero-midpoint} proves the claim.  At the level of
finite laws the proof retains the separate estimates
\[
 \mathbb E(D+E)\le2\varepsilon,\qquad
 \|\mathbb E s\|_2\le\sqrt{2\varepsilon},\qquad
 \|\mathbb E|s|\|_2\le\sqrt{6\varepsilon},\qquad
 \mathbb E_\pm\sqrt E\le2\sqrt\varepsilon.
\]
For the second assertion apply Theorem~\ref{stab:weighted} to
$z_\pm=h_\pm/R$.  Both individual deficits and the mixed deficit
are at most $d$, while $\eta=0$.  The conditional replacement
cost is therefore $\sqrt{2d}$ and its crossing cost is at most
$2\sqrt{6d}$.  Multiplication by $R$ proves the assertion, including
$d=0$.  Alternatively, projecting both endpoints onto increasing
finite ancestral sets containing $H$ preserves their sum and their
evaluations on $p$, and gives the same limit.
\end{proof}

\subsection{Keeping the original crossing coefficients}

The original coefficients give a second estimate in which segments
lying wholly outside the head remain visible as a separate error.

\begin{proposition}\label{stab:original-crossing}
Let $z_1,z_2\in B_Z$, $z=(z_1+z_2)/2$, and let $H$ be finite
ancestral.  For a head test $p$, put $\delta=1-z(p)$ and
$\kappa=\|(z(e_g))_{g\in G_H}\|_2$.  Then
\begin{equation}\label{stab:gate-original-bound}
 \|Q_Hz_i\|\le2\kappa+(6\sqrt2+6)\sqrt\delta.
\end{equation}
In particular, choosing a norming head test gives
\begin{equation}\label{stab:original-crossing-bound}
 \|Q_Hz_i\|
 \le2\|Q_Hz\|+(6\sqrt2+6)\sqrt{1-\|P_Hz\|}
 \le2\|Q_Hz\|+15\sqrt{1-\|P_Hz\|}.
\end{equation}
\end{proposition}

\begin{proof}
Use finite atom laws first, with equal mixture $\mathbb E$.
The decomposition $Q_HA=C(A)+O(A)$ from the toolkit separates
crossing and wholly exterior segments.  Let $c_g(A)$ and $s_g(A)$
be the original crossing coefficient and its head sum.  The
pointwise estimates \eqref{stab:coef:exterior} and the gate identity
\eqref{stab:coef:gate-coordinates} imply
\[
 \|\mathbb Ec\|_2\le\kappa+\sqrt{2\delta},\qquad
 \|\mathbb Es-\mathbb Ec\|_2\le\sqrt{2\delta}.
\]
By \eqref{stab:signed-unsigned},
$\|\mathbb E|s|\|_2\le\|\mathbb Es\|_2+2\sqrt\delta$.
The coordinate inequality $\big||c_g|-|s_g|\big|\le|c_g-s_g|$
and Jensen's inequality therefore give
\begin{equation}\label{stab:original-mass}
 \|\mathbb E|c|\|_2
 \le\kappa+(3\sqrt2+2)\sqrt\delta.
\end{equation}
Lemma~\ref{stab:coef:tail-assembly} bounds the crossing mean of
either conditional law by $2\|\mathbb E|c|\|_2$.  Its exterior
mean costs at most
\[
 \mathbb E_i\|O(A)\|
 \le\sqrt{2\mathbb E_iL(A)}\le2\sqrt\delta.
\]
Their sum proves \eqref{stab:gate-original-bound}.  For a norming
head test, $\delta=1-\|P_Hz\|$; gate contraction gives
$\kappa\le\|Q_Hz\|$, and $6\sqrt2+6<15$.

Keeping the weaker conditional exterior estimate
$\mathbb E_i\|O(A)\|\le2\sqrt{2\delta}$ instead records the
alternative budget
\begin{equation}\label{stab:sixteen-budget}
 \|Q_Hz_i\|\le2\kappa+(8\sqrt2+4)\sqrt\delta
 <2\kappa+16\sqrt\delta\qquad(\delta>0).
\end{equation}
For $\delta=0$ the last inequality is non-strict.  Either the
one-piece truncation or the two-piece comparison in
Lemma~\ref{stab:coef:head-conflict} supplies the required opposite-sign
estimate.  Thus this original-coefficient calculation is compatible
with both surgeries.  Normalized finite approximants and norm
continuity give the result for completed endpoints.
\end{proof}

\subsection{A small tail at an arbitrary center}

For a fixed center we can use an ordered comparison: positive
crossings from the first sample are matched only against negative
crossings from the second.  The deterministic step is the one-piece
truncation in Lemma~\ref{stab:coef:head-conflict}; its ordered
expectation keeps a different budget.

\begin{proposition}\label{stab:small-center-tail}
Let $0<e<1/4$, let $x\in Z$ satisfy $\|x\|=1$, and let $H$ be
finite ancestral with $\|Q_Hx\|\le e$.  Suppose $p=P_Hp$,
$\rho(p)=1$, and $x(p)\ge1-e$.  For every $u\in\ker P_H$,
\begin{equation}\label{stab:small-center-implication}
 \max\{\|x+u\|,\|x-u\|\}\le1+e
 \quad\Longrightarrow\quad \|u\|\le25\sqrt e.
\end{equation}
\end{proposition}

\begin{proof}
Fix a finite ancestral $K\supset H$ and represent
$P_K(x\pm u)/(1+e)$ by finite atom laws.  Under each law the
expectation on $p$ is at least $(1-e)/(1+e)\ge1-2e$.
The pointwise budget gives
\[
 \mathbb E_\pm E(A)\le4e,\qquad
 \mathbb E_\pm D(A)\le4e,\qquad
 \|\mathbb E_\pm(A-A^\circ)\|\le2\sqrt e.
\]
Let $l_g,n_g$ be the positive and negative crossing masses of the
equal mixture after replacement.  Its mean before replacement is
$P_Kx/(1+e)$, so gate contraction gives
$\|l-n\|_2\le2\sqrt e+e$.

Take independent outcomes of this mixed law.  At each gate compare
only a positive head sum in the first outcome with a negative head
sum in the second.  If the sums of the longer and shorter nested
heads are $b,b'$, truncation leaves the sum $b-b'$ and gains
$(b-b')^2-b^2=(b')^2-2bb'\ge2|bb'|$ in the squared test.
As in Lemma~\ref{stab:coef:head-conflict}, different gates require
different original segments, so the gains total at most $D(A)+D(A')$.
Independence now gives the one-orientation estimate
$2\sum_gl_gn_g\le\mathbb E(D(A)+D(A'))\le8e$.  Hence
\begin{equation}\label{stab:directional-mass}
 \sum_gl_gn_g\le4e,\qquad
 \|l+n\|_2\le\sqrt{(2\sqrt e+e)^2+16e}.
\end{equation}
Rooted-mass recovery bounds each conditional replaced tail by
$2\|l+n\|_2$.  Either proof in
Lemma~\ref{stab:coef:tail-assembly} applies; in its finite-test
proof all maximizing segments may be taken in the finite set $K$.

Using the sufficient bound $\|I-P_H\|\le2$, restriction of each
replacement error costs at most $4\sqrt e$.  Half the difference
of the conditional tails is $P_Ku/(1+e)$, since the center cancels.
Consequently
\[
 \|P_Ku\|\le(1+e)
 \left[2\sqrt{(2\sqrt e+e)^2+16e}+4\sqrt e\right]
 \le\frac54(\sqrt{89}+4)\sqrt e<25\sqrt e.
\]
Here $e\le\sqrt e/2$.  Coordinate exhaustion proves the claim.
\end{proof}

\subsection{From endpoint stability to averaged midpoint gain}

We use the averaged modulus of \eqref{eq:averaged-modulus}.
A norming-functional kernel makes both endpoint norms at least one;
a finite head alone gives this lower bound for a head-supported
center.  These two choices yield the following consequences.

\begin{corollary}\label{stab:averaged-gain}
For either coordinate model,
\[
 \widehat\delta_Z(t)\ge t^2/1024\qquad(0<t<1).
\]
For every $t>0$ and every $0<e<1/4$ with $25\sqrt e<t$, one also
has $\widehat\delta_Z(t)\ge e/2$.
\end{corollary}

\begin{proof}
For the first assertion set $d=t^2/256$.  Let $x_0$ be a finitely
supported unit vector and choose $H$ containing its support.
For $y\in\ker P_H$, $\|y\|\ge1$, both endpoint norms are at
least one.  If both were at most $1+d$, their normalized midpoint
$x_0/(1+d)$ would have norm at least $1-d$.
Equation~\eqref{stab:supported-bound} would give
\[
 \frac{t\|y\|}{1+d}\le(2+2\sqrt6)\sqrt d
 <7\sqrt d=7t/16<t/2,
\]
whereas $d<1$ makes the left side greater than $t/2$.
The average therefore exceeds $1+d/2$.  For an arbitrary unit
$x$, choose a finitely supported unit $x_0$ with
$\|x-x_0\|<d/4$ and use its head subspace.  The average loses
less than $d/4$, leaving gain $d/4=t^2/1024$.

For the second assertion choose a head with $\|Q_Hx\|\le e$ and
a norming head test with $x(p)\ge1-e$.  Intersect $\ker P_H$
with the kernel of a norming functional for $x$ on $Z$.
Both actual endpoints have norm at least one.  For $u=ty$,
$\|y\|\ge1$, Proposition~\ref{stab:small-center-tail} forces
one endpoint to exceed $1+e$, so their average exceeds $1+e/2$.
Taking the modulus proves the claim.
\end{proof}
